% P15-U1: the pairwise-independent and random-series strong-law proofs.
\chapter{Strong Laws of Large Numbers}
\label{sec:p15-strong-laws-of-large-numbers}

A strong law identifies the limiting average on one event of probability
one. For a common integrable distribution, pairwise independence already
suffices. A different proof for mutually independent variables uses
convergence of a random series and gives a useful weighted criterion for
nonidentical distributions.

Both source presentations are retained: Durrett's author Version 5 dated
11 January 2019, Theorem 2.4.1 and its Lemmas 2.4.2--2.4.4
\cite[pp.~76--78]{DurrettPTEAuthor2019}, and Kronecker's lemma and
the second strong-law proof, Theorems 2.5.9--2.5.10
\cite[pp.~85--87]{DurrettPTEAuthor2019}. The first proof uses pairwise
variance addition, geometric subsequences and monotonicity. The second
uses the mutually independent random-series theorem already proved in
Chapter~\ref{sec:p15-independent-random-series}.

\section{Averages and the truncation estimates}
\label{sec:p15-strong-law-truncation}

\begin{hbdefinition}{def:p15-strong-law-averages}{Sample averages and an integrable common law}
For real random variables \(X_1,X_2,\ldots\) on
\((\Omega,\mathcal F,\mathbb P)\), write
\[
 S_0=0,\qquad S_n=\sum_{k=1}^nX_k,\qquad
 \overline X_n=S_n/n\quad(n\geq1).
\]
A strong law with finite value \(\mu\in\mathbb R\) is the
assertion \(\mathbb P(\overline X_n\to\mu)=1\).
An integrable common law means that the \(X_n\) are identically
distributed and \(\mathbb E|X_1|<\infty\); then
\(\mu=\mathbb EX_1=\mathbb EX_n\) is finite for every \(n\).
Pairwise independence means independence of every two distinct variables.
The abbreviation i.i.d. requires mutual independence of every finite
subfamily as well as identical distributions, as in
Definition~\ref{def:p15-independent-random-elements}.

For each integer \(k\geq1\), define the moving deletion truncation
and its partial sums by
\[
 Y_k=X_k\mathbf1_{\{|X_k|\leq k\}},\qquad
 T_0=0,\qquad T_n=\sum_{k=1}^nY_k.
\]
The cutoff is the index \(k\); observations of magnitude exactly
\(k\) are retained. Different \(Y_k\) need not have identical
distributions, even when the \(X_k\) do.
\index{strong law of large numbers}\index{sample average}
\end{hbdefinition}

\begin{hblemma}{thm:p15-cesaro-average}{proved=annex:p15.cesaro-average}{Averaging a convergent numerical sequence}
If a real sequence \(u_k\) converges to \(u\in\mathbb R\), then
\[
 \frac1n\sum_{k=1}^nu_k\longrightarrow u.
\]
In particular this applies to a sequence of finite expectations whose
individual values converge. No independence is involved.
\end{hblemma}

\begin{hblemma}{thm:p15-strong-law-tail-bounds}{proved=annex:p15.strong-law-tail-bounds}{Tail integrals and the reciprocal-square bound}
For a nonnegative real random variable \(W\), possibly with infinite
moments,
\[
 \mathbb EW=\int_0^\infty\mathbb P(W>t)\,dt,
 \qquad
 \mathbb EW^2=\int_0^\infty2t\mathbb P(W>t)\,dt.
\]
The integrals are with respect to Borel Lebesgue measure. Moreover,
\[
 \sum_{k\geq1}\mathbb P(W>k)\leq\mathbb EW,
 \qquad
 2y\sum_{k>y}k^{-2}\leq4\quad(y\geq0).
\]
The second sum ranges over positive integer \(k\).
\end{hblemma}

\begin{hbtheorem}{thm:p15-strong-law-truncation}{proved=annex:p15.strong-law-truncation}{All truncation estimates from a common first moment}
Suppose only that the \(X_k\) have an integrable common law.
Independence is not required for the following conclusions:
\[
 \begin{aligned}
 \sum_{k\geq1}\mathbb P(X_k\ne Y_k)
       &=\sum_{k\geq1}\mathbb P(|X_1|>k)
                      \leq\mathbb E|X_1|<\infty,\\
 \sum_{k\geq1}\frac{\operatorname{Var}(Y_k)}{k^2}
       &\leq\sum_{k\geq1}\frac{\mathbb EY_k^2}{k^2}
                      \leq4\mathbb E|X_1|<\infty.
 \end{aligned}
\]
Almost surely, \(X_k=Y_k\) for all sufficiently large \(k\).
On that event there is a finite bound \(R(\omega)\) with
\(|S_n(\omega)-T_n(\omega)|\leq R(\omega)\) for every \(n\).
Consequently \((S_n-T_n)/n\to0\) almost surely. Also,
\[
 \mathbb EY_k\longrightarrow\mu,
 \qquad
 \frac{\mathbb ET_n}{n}\longrightarrow\mu.
\]
Pairwise independence of the \(X_k\), when supplied, is preserved
by these separate measurable truncations; mutual independence is
preserved when that stronger premise is supplied.
\end{hbtheorem}

\section{The pairwise independent strong law}
\label{sec:p15-pairwise-strong-law}

\begin{hblemma}{thm:p15-geometric-subsequence}{proved=annex:p15.geometric-subsequence}{Geometric indices and summable subsequence deviations}
Let \(\alpha>1\), and for integers \(j\geq1\) set
\(k(j)=\lfloor\alpha^j\rfloor\), where \(\lfloor t\rfloor\)
is the greatest integer at most \(t\). Then \(k(j)\geq1\),
\(k(j)\to\infty\), and \(k(j+1)/k(j)\to\alpha\).
For every positive integer \(m\),
\[
 \sum_{j:k(j)\geq m}k(j)^{-2}
                     \leq\frac{4}{1-\alpha^{-2}}\,m^{-2}.
\]
If \(Y_m\) are pairwise independent with finite second moments and
\(\sum_m\operatorname{Var}(Y_m)/m^2<\infty\), put
\(T_\ell=\sum_{m=1}^\ell Y_m\). Then for every
\(\varepsilon>0\),
\[
 \sum_{j\geq1}
 \mathbb P\!\left(|T_{k(j)}-\mathbb ET_{k(j)}|>
                                       \varepsilon k(j)\right)
 \leq\frac{4\varepsilon^{-2}}{1-\alpha^{-2}}
           \sum_{m\geq1}\frac{\operatorname{Var}(Y_m)}{m^2}<\infty.
\]
In particular,
\((T_{k(j)}-\mathbb ET_{k(j)})/k(j)\to0\) almost surely.
\end{hblemma}

\begin{hbtheorem}{thm:p15-etemadi-strong-law}{proved=annex:p15.etemadi-strong-law}{The strong law under pairwise independence}
Let \(X_1,X_2,\ldots\) be pairwise independent and identically
distributed real random variables with \(\mathbb E|X_1|<\infty\).
Then
\[
 \frac1n\sum_{k=1}^nX_k\longrightarrow\mathbb EX_1
                                      \qquad\text{almost surely}.
\]
This is the full statement of Durrett's Theorem 2.4.1, whose proof is
attributed there to Etemadi. The annex retains the source's nonnegative
reduction, geometric-subsequence argument and interpolation, with
countably many thresholds and ratios furnishing a common event of
probability one.
\end{hbtheorem}

\section{Kronecker's lemma and the independent-series route}
\label{sec:p15-kronecker-strong-law}

\begin{hbdefinition}{def:p15-kronecker-weights}{Weights and weighted series}
A positive divergent nondecreasing weight sequence is a sequence of
finite real numbers satisfying
\[
 0<a_1\leq a_2\leq\cdots,\qquad a_n\longrightarrow\infty.
\]
Repeated weights are allowed. For real numbers \(x_n\), the weighted
series is \(\sum_{n\geq1}x_n/a_n\); convergence means a finite
limit of its partial sums and need not be absolute. For the summation
identity below, set \(a_0=0\), \(b_0=0\) and
\(b_n=\sum_{k=1}^nx_k/a_k\) for \(n\geq1\). No term involving
division by \(a_0\) is defined or used.
\index{Kronecker's lemma}\index{weighted series}
\end{hbdefinition}

\begin{hbtheorem}{thm:p15-kronecker-lemma}{proved=annex:p15.kronecker-lemma}{Kronecker's lemma with the weight domain explicit}
For the positive divergent nondecreasing weights above, convergence of
\(\sum_nx_n/a_n\) implies
\[
 \frac1{a_n}\sum_{k=1}^nx_k\longrightarrow0.
\]
The proof uses the exact summation-by-parts identity
\[
 \frac1{a_n}\sum_{k=1}^nx_k
   =b_n-\sum_{m=1}^n\frac{a_m-a_{m-1}}{a_n}\,b_{m-1}.
\]
More generally, the conclusion still holds if every \(a_n\ne0\),
\(a_n\to\infty\), and the weights are positive and nondecreasing
from some index onward. Thus any finite initial nonzero weights can be
removed before applying the positive-weight formulation.
\end{hbtheorem}

\begin{hbcorollary}{thm:p15-weighted-independent-strong-law}{proved=annex:p15.weighted-independent-strong-law}{A weighted criterion without identical distributions}
Let \(X_n\) be mutually independent real random variables with finite
second moments and means \(\mu_n=\mathbb EX_n\). Let \(a_n\)
be positive divergent nondecreasing weights. If
\[
 \sum_{n\geq1}\frac{\operatorname{Var}(X_n)}{a_n^2}<\infty,
\]
then
\[
 \frac1{a_n}\sum_{k=1}^n(X_k-\mu_k)\longrightarrow0
                                      \qquad\text{almost surely}.
\]
In particular, if \(a_n=n\) and
\(n^{-1}\sum_{k=1}^n\mu_k\to\mu\in\mathbb R\), then
\(\overline X_n\to\mu\) almost surely. Identical distributions
are not a premise of this criterion.
\end{hbcorollary}

\begin{hbtheorem}{thm:p15-iid-series-strong-law}{proved=annex:p15.iid-series-strong-law}{The integrable i.i.d. strong law by random series}
Let \(X_1,X_2,\ldots\) be i.i.d. real random variables with
\(\mathbb E|X_1|<\infty\) and \(\mu=\mathbb EX_1\). Then
\[
 \overline X_n\longrightarrow\mu\qquad\text{almost surely}.
\]
This is Durrett's Theorem 2.5.10 with its distinct proof: the moving
truncations satisfy the common-law bounds, the series
\(\sum_k(Y_k-\mathbb EY_k)/k\) converges almost surely, Kronecker's
lemma controls the centered averages, and the truncated means tend to
\(\mu\). The proof does not invoke
Theorem~\ref{thm:p15-etemadi-strong-law}. Mutual independence is used
to apply the random-series theorem;
Theorem~\ref{thm:p15-etemadi-strong-law} establishes the stronger
pairwise-independent conclusion by a different argument.
\end{hbtheorem}

\section{Applications and boundaries}
\label{sec:p15-strong-law-applications}

\begin{hbtheorem}{thm:p15-strong-law-applications}{proved=annex:p15.strong-law-applications}{Integrable observations, frequencies and two computed averages}
Let \(X_n\) be pairwise independent with common real law \(\nu\).
For any fixed Borel function \(g:\mathbb R\to\mathbb R\) with
\(\int|g|\,d\nu<\infty\),
\[
 \frac1n\sum_{k=1}^ng(X_k)\longrightarrow\int g\,d\nu
                                           \quad\text{almost surely}.
\]
For any specified countable collection of such functions, all these
limits hold on a common event of probability one. In particular the
observed frequency of a fixed Borel event \(B\) tends to \(\nu(B)\).
Two explicit calculations are:
\begin{enumerate}
\item If \(U_k\) are i.i.d. uniform on \([0,1]\), then
\(n^{-1}\sum_{k=1}^nU_k^2\to\int_0^1u^2\,du=1/3\) almost surely.
\item If \(X_k\) are i.i.d. with density
\(f(x)=(3/2)x^{-5/2}\mathbf1_{[1,\infty)}(x)\), then
\(\overline X_n\to3\) almost surely, although
\(\mathbb EX_1^2=\infty\).
\end{enumerate}
The second example shows that a finite variance is unnecessary for the
integrable strong laws. The frequency conclusion is for fixed events
or a specified countable collection; uniform convergence over all real
thresholds is a separate assertion.
\end{hbtheorem}

\begin{hbtheorem}{thm:p15-pairwise-block-example}{proved=annex:p15.pairwise-block-example}{A strong law beyond mutual independence}
Let \(U_j,V_j\), \(j\geq1\), together be mutually independent
fair \(\{-1,1\}\)-valued signs. Define
\[
 X_{3j-2}=U_j,\qquad X_{3j-1}=V_j,\qquad X_{3j}=U_jV_j.
\]
The \(X_n\) have identical fair-sign laws and are pairwise independent,
but they are not mutually independent. Nevertheless
\(n^{-1}\sum_{k=1}^nX_k\to0\) almost surely by the pairwise strong
law. Each triple has product one, and its probability of three positive
signs is \(1/4\), rather than the product \(1/8\) of its marginals.
\end{hbtheorem}

\begin{hbtheorem}{thm:p15-strong-law-counterexamples}{proved=annex:p15.strong-law-counterexamples}{Why the law and weight hypotheses matter}
The following are fully specified boundaries.
\begin{enumerate}
\item If \(X_n=\varepsilon\) for all \(n\), with one fair
\(\{-1,1\}\)-valued sign \(\varepsilon\), the variables have an
identical bounded law and mean zero, but
\(\overline X_n=\varepsilon\) at every outcome. Pairwise independence
has failed.
\item Let independent \(X_n\) have mass \(1/(2n)\) at each of
\(n\) and \(-n\), and mass \(1-1/n\) at zero. Each mean is zero
and \(\mathbb E|X_n|=1\), but the averages have no finite limit
almost surely. Thus even a uniform bound on the individual first
absolute moments cannot replace the common-law hypothesis by itself.
\item Let i.i.d. positive \(X_n\) satisfy
\(\mathbb P(X_n=2^j)=2^{-j}\) for integers \(j\geq1\).
Their first moments are infinite, and
\(\overline X_n\to\infty\) almost surely. The finite-mean conclusion
of the integrable theorem cannot be asserted in this regime.
\item For Kronecker's lemma, take
\(a_{2j-1}=j^3\), \(a_{2j}=j\), \(x_{2j-1}=j\),
\(x_{2j}=-1/j\). The weights are positive and tend to infinity,
and \(\sum_n|x_n/a_n|<\infty\), but
\(a_{2j}^{-1}\sum_{k=1}^{2j}x_k\to\infty\).
The weights are not eventually nondecreasing.
Also, constant weights \(a_n=1\), with \(x_1=1\) and all later
\(x_n=0\), show why divergence of the weights is required.
\end{enumerate}
The independent sequences with these discrete laws exist by
Theorem~\ref{thm:p15-independent-sequence-existence}.
\end{hbtheorem}

\begin{proofboundary}{def:p15-strong-law-next-boundary}{Infinite means, rates and uniform empirical convergence}
Both integrable source strong laws and Kronecker's lemma are proved
here with their complete proof routes. The remaining general
infinite-mean law is Durrett's Theorem 2.4.5, allowing an infinite
positive-part expectation and a finite negative-part expectation;
the positive dyadic example above establishes only its specified
special case. The complete renewal and uniform empirical-distribution
applications, Theorems 2.4.7 and 2.4.9, remain separate constructions.
The next random-series results are the finite-variance rate and the
Marcinkiewicz--Zygmund rate, Theorems 2.5.11--2.5.12, followed by the
infinite-mean normalization theorem 2.5.13
\cite[pp.~78--80, 87--90]{DurrettPTEAuthor2019}. No convergence rate,
uniform empirical limit or general infinite-mean normalization is
inferred merely from the integrable laws above.
\end{proofboundary}
